\documentclass[10pt,a4paper]{amsart}
\usepackage[margin=19mm]{geometry}
\usepackage{amsmath,amssymb,mathtools}
\usepackage[scr=rsfs,cal=euler]{mathalfa}
\usepackage{xfrac}
\usepackage[unicode,hidelinks,bookmarks=false]{hyperref}
\newcommand{\ie}{i.e.\ }
\newcommand{\cf}{cf.\ }
\newcommand{\resp}{respectively\ }
\newcommand{\Subsection}{\subsection}
\providecommand{\nobreakdash}{\nobreak\hbox{-}}
\setlength{\emergencystretch}{2em}
\multlinegap0pt
\newtheorem{lemma}{Lemma}
\newcommand{\N}{\mathbb{N}}
\newcommand{\R}{\mathbb{R}}
\newcommand{\calK}{\mathcal{K}}
\newcommand{\calW}{\mathcal{W}}
\newcommand{\calX}{\mathcal{X}}
\DeclareMathOperator{\clip}{clip}
\newcommand\defeq{\vcentcolon=}
\renewcommand{\ge}{\geqslant}
\renewcommand{\le}{\leqslant}
\DeclareMathOperator{\solid}{solid}
\title{Norm Bounds for Sparse Random Tensors and\\ Spectral Gap of Random Hypergraphs}
\author{Kevin Lucca, Lucas Pesenti}
\date{Source: arXiv:2607.07308v1 (CC BY 4.0)}
\begin{document}
\maketitle

\noindent\emph{Source and licence.} Norm Bounds for Sparse Random Tensors and\\ Spectral Gap of Random Hypergraphs. Version arXiv:2607.07308v1 (CC BY 4.0), distributed by arXiv under the Creative Commons Attribution 4.0 International licence, http://creativecommons.org/licenses/by/4.0/, as recorded in the arXiv licence metadata for this exact version and shown on the abstract page https://arxiv.org/abs/2607.07308v1. Full licence notice: \texttt{licenses/arxiv-2607.07308-CC-BY-4.0.txt}.\par\smallskip

\noindent\textit{Editorial scope.} Counted unit: the article's lemma lem:solidContain (source0.tex lines 798--803). The article's complete original proof of it is delivered with it (source0.tex lines 805--857, one \texttt{proof} environment of 53 lines). No mathematics is written, repaired or re-derived: the statement and the proof are the article's own lines, in the article's own notation, and no retained line is substituted or re-flowed.  No further source-local context is needed: every cross-reference of every retained line resolves inside the two delivered spans.  Nothing was omitted from any retained span, so the package carries no omission record.  No retained line contains a citation, so the wrapper carries no bibliography and no bibliography entry.  No retained line contains a figure environment, an image-inclusion command or drawing code, so no image is delivered and the package declares no asset dependency.  Editorial formatting only, in addition: an AMS \texttt{amsart} preamble that restates the article's own declarations and macros used by the retained lines, namely the environments \texttt{lemma} on separate counters, together with the standard packages those lines need.  The retained environments print in this document as Lemma 1 in order of appearance, whereas the article prints the same items with the numbering of its own full document.  The complete native source of the article as distributed in the arXiv e-print for this exact version is bundled unchanged as \texttt{sources/204694/source0.tex}.
\begin{lemma} \label{lem:solidContain}
    We have
    \[  
        \calX_P \subseteq \solid\left(\bigcup\limits_{j=0}^r \calW_j\right)\,.
    \]
\end{lemma}

\begin{proof}
    It suffices to consider
    vectors with non-negative entries. Let $x_1, \ldots, x_r\in \R^n_{\ge 0}$ be unit vectors.
    Given
    $k=(k_1, \ldots, k_r)\in\N^r$ and $i\in [r]$, write
    $A_{k,i} \defeq \{j\in [n] : (x_i)_j\in [r^{-k_i-1}, r^{-k_i}]\}$.
    Since each $x_i$ is a unit vector,
    \begin{equation}
        \sum_{k\ge 0} |A_{k,i}| r^{-2k_i-2} \le 1\implies |A_{k,i}|\le r^{2k_i+2}\,.\label{eq:counting}
    \end{equation}
    Taking tensor products of this level set
    decomposition gives
    \begin{equation}
        \bigotimes_{i\in [r]} x_i - \clip_\tau\left(\bigotimes_{i\in [r]} x_i\right)\le \sum_{\substack{k_1, \ldots, k_r\ge 0\\\prod\limits_{i\in [r]}r^{-k_i}>\tau}} \bigotimes_{i\in [r]} \frac{\mathbf 1_{A_{k,i}}}{r^{k_i}}\,.\label{eq:mainDiscret}
    \end{equation}
    We cover the indices of the sum on the right-hand side of~\eqref{eq:mainDiscret} by
    \begin{align*}
        \calK_{0} &= \left\{(k_1, \ldots, k_r)\in \N^r:
        \frac 1r \sum_{i=1}^r |A_{k,i}|\le \prod_{i=1}^r \phi(A_{k,i})\text{ and } \prod_{i=1}^r |A_{k,i}|\le r^{2r} n^2\log^2 n\right\}\,,\\
        \calK_{j} &= \left\{(k_1, \ldots, k_r)\in \N^r : \frac 1r \sum_{i=1}^r |A_{k,i}| > \prod_{i=1}^r \phi(A_{k,i}) \text{ and } |A_{k,j}| = \max_{i\in [r]} |A_{k,i}|\right\}\quad \forall j\in [r]\,.
    \end{align*}
    First, we claim that 
    the sets $\calK_0,\ldots,\calK_r$
    cover every index appearing in the sum on 
    the right-hand side
    of~\eqref{eq:mainDiscret}. Indeed, by~\eqref{eq:counting}, we have
    \[
        \prod_{i\in [r]} r^{-k_i}>\tau \implies \prod_{i=1}^r |A_{k,i}|\le r^{2r} \prod_{i=1}^r r^{2k_i}\le \frac {r^{2r}} {\tau^2} = r^{2r} n^2 \log^2 n\,,
    \]
    so the second condition
    in the definition of $\calK_0$
    is automatically satisfied.

    Second, we show that, for each $j\in \{0, \ldots, r\}$, the sum over $\calK_j$
    lies in the solid convex hull
    of the corresponding $\calW_j$.
    
    \begin{lemma}\label{lemma:bSum}
    We have
    \[
        \sum_{k\in \calK_0} \bigotimes_{i\in [r]} \frac {\mathbf 1_{A_{k,i}}} {r^{k_i}} \in \frac12 \solid(\calW_0)\,.
    \]
    \end{lemma}
    
    \begin{lemma}\label{lemma:uSum}
        For any $j\in [r]$, we have
        \[\sum_{k\in \calK_j} \bigotimes_{i\in [r]} \frac {\mathbf 1_{A_{k,i}}} {r^{k_i}}  \in \frac 1{2r}  \solid(\calW_j)\,.
        \]
    \end{lemma}
    
    \noindent Lemmas~\ref{lemma:bSum}
    and~\ref{lemma:uSum} complete the proof of Lemma~\ref{lem:solidContain}.
\end{proof}

\end{document}
